\documentclass[11pt,a4paper]{article}

% =============================================================================
% ENCODING
% =============================================================================
\usepackage[utf8]{inputenc}
\usepackage[T1]{fontenc}

% =============================================================================
% PACKAGES
% =============================================================================
\usepackage{amsmath,amssymb,amsthm}
\usepackage{mathrsfs}
\usepackage[margin=1in]{geometry}
\usepackage{enumitem}
\usepackage[backend=biber]{biblatex}
\addbibresource{../release/references.bib}
\usepackage[unicode,hidelinks]{hyperref}
\hypersetup{
  pdftitle={Contractive Projection Dynamics on Hilbert Spaces},
  pdfauthor={Johnny Andrey P\'erez Ram\'irez},
  pdfsubject={Rigid Identity Framework -- Volume 0: Mathematical Foundation},
  pdfkeywords={Hilbert space, contractive dynamics, metric projection, fixed point, attractor, non-collapsibility, spectral gap}
}

% Pagination
\raggedbottom
\widowpenalty=10000
\clubpenalty=10000

% =============================================================================
% THEOREM ENVIRONMENTS
% =============================================================================
\theoremstyle{plain}
\newtheorem{theorem}{Theorem}[section]
\newtheorem{lemma}[theorem]{Lemma}
\newtheorem{proposition}[theorem]{Proposition}
\newtheorem{corollary}[theorem]{Corollary}

\theoremstyle{definition}
\newtheorem{definition}[theorem]{Definition}
\newtheorem{axiom}[theorem]{Axiom}
\newtheorem{hypothesis}[theorem]{Hypothesis}
\newtheorem{remark}[theorem]{Remark}
\newcommand{\Hilb}{\mathcal{H}}
\newcommand{\I}{\mathcal{I}}
\newcommand{\U}{\mathcal{U}}
\newcommand{\A}{\mathcal{A}}
\newcommand{\norm}[1]{\left\lVert #1 \right\rVert}
\newcommand{\inner}[2]{\left\langle #1, #2 \right\rangle}

% ... (omitted lines)

\title{%
\textbf{Contractive Projection Dynamics on Hilbert Spaces}\\[0.5em]
\large Foundation for the Rigid Identity Framework
}

\author{Johnny Andrey P\'{e}rez Ram\'{i}rez}

\date{}

% =============================================================================
% DOCUMENT
% =============================================================================
\begin{document}

\maketitle

% =============================================================================
% ABSTRACT
% =============================================================================
\begin{abstract}
We establish a canonical mathematical foundation for contractive dynamical systems on Hilbert spaces with metric projections onto closed convex subsets. We prove existence and uniqueness of projections, contractivity of the induced transition operator, existence of unique attractors, global compactness, and structural non-collapsibility. All results are derived from four minimal hypotheses (H1--H4) using standard functional analysis. This core serves as the rigorous foundation for the Rigid Identity Framework, where \textit{Rigid Identity} and \textit{Phenomenological Regimes} become direct corollaries and the mathematical dependencies used by \textit{Null-Regime Instability} are made explicit. No ontological, phenomenological, or semantic interpretations are assumed; all results are purely structural.
\end{abstract}

\noindent\textbf{Release governance note.}\enspace Methods, governance conventions, claim boundaries, and terminology policy for this release are maintained in the project's canonical governance hub (\textit{Methods Governance Hub}, release version~1). This document adheres to that hub as its editorial and methodological reference.

\tableofcontents
\newpage

% =============================================================================
% SECTION 1: INTRODUCTION
% =============================================================================
\section{Introduction}

\subsection{Purpose}

This document establishes the \textbf{canonical mathematical core} of the Rigid Identity Framework. All subsequent results in \textit{Rigid Identity}, \textit{Phenomenological Regimes}, and \textit{Null-Regime Instability} are derived as corollaries of the theorems proven here.

\subsection{Methodology}

We employ standard techniques from:
\begin{itemize}
    \item Functional analysis on Hilbert spaces
    \item Theory of metric projections
    \item Banach fixed-point theorem
    \item Compactness arguments
\end{itemize}

\subsection{Scope}

This is \textbf{pure mathematics}. No interpretation is required or assumed. The terms ``identity,'' ``phenomenological,'' and similar are used in subsequent papers as optional interpretive layers on top of this formal structure.

% =============================================================================
% SECTION 2: HYPOTHESES
% =============================================================================
\section{Minimal Hypotheses}

We establish four hypotheses that are both necessary and sufficient for all main results.

\begin{hypothesis}[H1: Projection Structure]\label{hyp:H1}
Let $\Hilb$ be a real Hilbert space with inner product $\inner{\cdot}{\cdot}$ and induced norm $\norm{\cdot}$. Let $\I \subset \Hilb$ be a non-empty, closed, and convex subset.
\end{hypothesis}

\begin{hypothesis}[H2: Contraction Operator]\label{hyp:H2}
Let $K: \Hilb \to \Hilb$ be a bounded linear operator with operator norm $\norm{K} < 1$.
\end{hypothesis}

\begin{hypothesis}[H3: Completeness]\label{hyp:H3}
The space $\Hilb$ is complete (as a Hilbert space, this is automatic).
\end{hypothesis}

\begin{hypothesis}[H4: Parameter Compactness]\label{hyp:H4}
Let $\U$ be a compact metric space (the parameter space). Let $\Gamma: \U \to \Hilb$ be a uniformly continuous function.
\end{hypothesis}

\subsection{Necessity and Sufficiency}

\begin{proposition}[Minimal Hypotheses]\label{prop:minimal}
Hypotheses H1--H4 are:
\begin{enumerate}
    \item \textbf{Necessary:} Removal of any hypothesis leads to failure of at least one main theorem.
    \item \textbf{Sufficient:} Together they imply all main theorems.
\end{enumerate}
\end{proposition}

\begin{proof}
Necessity is established by counterexamples in Appendix A. Sufficiency is the content of Theorems \ref{thm:projection}--\ref{thm:noncollapse}.
\end{proof}

% =============================================================================
% SECTION 3: PROJECTION OPERATOR
% =============================================================================
\section{The Projection Operator}

\begin{definition}[Metric Projection]\label{def:projection}
For $\Psi \in \Hilb$, define the metric projection onto $\I$ as:
\[
\aleph(\Psi) := \arg\min_{\Phi \in \I} \norm{\Psi - \Phi}
\]
\end{definition}

\begin{theorem}[Existence and Uniqueness of Projection]\label{thm:projection}
Under Hypothesis H1, for every $\Psi \in \Hilb$, the projection $\aleph(\Psi)$ exists and is unique.
\end{theorem}

\begin{proof}
\textbf{Existence.} Since $\I$ is non-empty and closed, and $\Hilb$ is complete, the infimum $d := \inf_{\Phi \in \I} \norm{\Psi - \Phi}$ is attained. Let $\{\Phi_n\}$ be a minimizing sequence with $\norm{\Psi - \Phi_n} \to d$. By the parallelogram law:
\[
\norm{\Phi_m - \Phi_n}^2 = 2\norm{\Psi - \Phi_m}^2 + 2\norm{\Psi - \Phi_n}^2 - 4\norm{\Psi - \frac{\Phi_m + \Phi_n}{2}}^2
\]
Since $\I$ is convex, $\frac{\Phi_m + \Phi_n}{2} \in \I$, so $\norm{\Psi - \frac{\Phi_m + \Phi_n}{2}} \geq d$. Thus:
\[
\norm{\Phi_m - \Phi_n}^2 \leq 2\norm{\Psi - \Phi_m}^2 + 2\norm{\Psi - \Phi_n}^2 - 4d^2 \to 0
\]
as $m, n \to \infty$. Hence $\{\Phi_n\}$ is Cauchy. By completeness, $\Phi_n \to \Phi^* \in \Hilb$. Since $\I$ is closed, $\Phi^* \in \I$. By continuity of the norm, $\norm{\Psi - \Phi^*} = d$, so the minimum is attained.

\textbf{Uniqueness.} Suppose $\Phi_1, \Phi_2 \in \I$ both minimize. Then by the parallelogram law:
\[
\norm{\Phi_1 - \Phi_2}^2 = 2\norm{\Psi - \Phi_1}^2 + 2\norm{\Psi - \Phi_2}^2 - 4\norm{\Psi - \frac{\Phi_1 + \Phi_2}{2}}^2
\]
\[
= 2d^2 + 2d^2 - 4\norm{\Psi - \frac{\Phi_1 + \Phi_2}{2}}^2 \leq 4d^2 - 4d^2 = 0
\]
since $\frac{\Phi_1 + \Phi_2}{2} \in \I$ by convexity. Thus $\Phi_1 = \Phi_2$.
\end{proof}

\begin{lemma}[Non-Expansiveness]\label{lem:nonexp}
The projection $\aleph: \Hilb \to \I$ is non-expansive:
\[
\norm{\aleph(\Psi_1) - \aleph(\Psi_2)} \leq \norm{\Psi_1 - \Psi_2}
\]
for all $\Psi_1, \Psi_2 \in \Hilb$.
\end{lemma}

\begin{proof}
Let $\Phi_i = \aleph(\Psi_i)$ for $i = 1, 2$. By the characterization of projections in Hilbert spaces, for all $\Phi \in \I$:
\[
\inner{\Psi_i - \Phi_i}{\Phi - \Phi_i} \leq 0
\]
Setting $\Phi = \Phi_2$ in the first inequality and $\Phi = \Phi_1$ in the second:
\[
\inner{\Psi_1 - \Phi_1}{\Phi_2 - \Phi_1} \leq 0, \quad \inner{\Psi_2 - \Phi_2}{\Phi_1 - \Phi_2} \leq 0
\]
Adding:
\[
\inner{\Psi_1 - \Phi_1 - (\Psi_2 - \Phi_2)}{\Phi_2 - \Phi_1} \leq 0
\]
\[
\inner{(\Psi_1 - \Psi_2) - (\Phi_1 - \Phi_2)}{\Phi_2 - \Phi_1} \leq 0
\]
\[
-\inner{\Psi_1 - \Psi_2}{\Phi_1 - \Phi_2} + \norm{\Phi_1 - \Phi_2}^2 \leq 0
\]
\[
\norm{\Phi_1 - \Phi_2}^2 \leq \inner{\Psi_1 - \Psi_2}{\Phi_1 - \Phi_2} \leq \norm{\Psi_1 - \Psi_2} \norm{\Phi_1 - \Phi_2}
\]
Dividing by $\norm{\Phi_1 - \Phi_2}$ (if nonzero):
\[
\norm{\Phi_1 - \Phi_2} \leq \norm{\Psi_1 - \Psi_2}
\]
\end{proof}

% =============================================================================
% SECTION 4: TRANSITION OPERATOR
% =============================================================================
\section{The Transition Operator}

\begin{definition}[Transition Operator]\label{def:transition}
For $u \in \U$, define the transition operator $T_u: \Hilb \to \Hilb$ by:
\[
T_u(\Psi) := \aleph(K\Psi + \Gamma(u))
\]
\end{definition}

\begin{theorem}[Contractivity]\label{thm:contraction}
Under Hypotheses H1--H2, for each $u \in \U$, $T_u$ is a strict contraction with Lipschitz constant $\norm{K} < 1$:
\[
\norm{T_u(\Psi_1) - T_u(\Psi_2)} \leq \norm{K} \cdot \norm{\Psi_1 - \Psi_2}
\]
\end{theorem}

\begin{proof}
For any $\Psi_1, \Psi_2 \in \Hilb$:
\begin{align*}
\norm{T_u(\Psi_1) - T_u(\Psi_2)} &= \norm{\aleph(K\Psi_1 + \Gamma(u)) - \aleph(K\Psi_2 + \Gamma(u))} \\
&\leq \norm{K\Psi_1 + \Gamma(u) - K\Psi_2 - \Gamma(u)} \quad \text{(Lemma \ref{lem:nonexp})} \\
&= \norm{K(\Psi_1 - \Psi_2)} \\
&\leq \norm{K} \cdot \norm{\Psi_1 - \Psi_2}
\end{align*}
Since $\norm{K} < 1$ by H2, $T_u$ is a strict contraction.
\end{proof}

% =============================================================================
% SECTION 5: ATTRACTOR EXISTENCE
% =============================================================================
\section{Existence of Unique Attractors}

\begin{theorem}[Unique Fixed Point]\label{thm:fixedpoint}
Under Hypotheses H1--H3, for each $u \in \U$, there exists a unique fixed point $f^*_u \in \Hilb$ such that:
\[
T_u(f^*_u) = f^*_u
\]
Moreover, for any $\Psi_0 \in \Hilb$:
\[
\lim_{n \to \infty} T_u^n(\Psi_0) = f^*_u
\]
with exponential convergence rate $\norm{K}^n$.
\end{theorem}

\begin{proof}
By Theorem \ref{thm:contraction}, $T_u$ is a strict contraction on the complete metric space $(\Hilb, \norm{\cdot})$. By the Banach Fixed-Point Theorem, there exists a unique fixed point $f^*_u$ and convergence holds with rate:
\[
\norm{T_u^n(\Psi_0) - f^*_u} \leq \frac{\norm{K}^n}{1 - \norm{K}} \norm{T_u(\Psi_0) - \Psi_0}
\]
\end{proof}

\begin{definition}[Global Attractor]\label{def:attractor}
The global attractor is:
\[
\A^* := \{f^*_u : u \in \U\}
\]
\end{definition}

\begin{theorem}[Compactness of Attractor]\label{thm:compactness}
Under Hypotheses H1--H4, the global attractor $\A^*$ is compact in $\Hilb$.
\end{theorem}

\begin{proof}
Define $F: \U \to \Hilb$ by $F(u) = f^*_u$. We show $F$ is continuous.

Let $u_n \to u$ in $\U$. We have:
\[
T_{u_n}(f^*_{u_n}) = f^*_{u_n}, \quad T_u(f^*_u) = f^*_u
\]

Consider any fixed $\Psi_0 \in \Hilb$. Then:
\[
f^*_{u_n} = \lim_{k \to \infty} T_{u_n}^k(\Psi_0), \quad f^*_u = \lim_{k \to \infty} T_u^k(\Psi_0)
\]

By uniform continuity of $\Gamma$ (H4), $\Gamma(u_n) \to \Gamma(u)$. For any fixed $k$:
\[
\norm{T_{u_n}^k(\Psi_0) - T_u^k(\Psi_0)} \to 0 \text{ as } n \to \infty
\]
by induction on $k$ and continuity of $\aleph$.

Taking $k \to \infty$ and $n \to \infty$ appropriately, we get $f^*_{u_n} \to f^*_u$.

Since $\U$ is compact and $F$ is continuous, $\A^* = F(\U)$ is compact as the continuous image of a compact set.
\end{proof}

% =============================================================================
% SECTION 6: NON-COLLAPSIBILITY
% =============================================================================
\section{Non-Collapsibility}

\begin{definition}[Constant Functions]
Let $\mathcal{N} := \{f \in \Hilb : f \text{ is constant}\}$ denote the subspace of constant elements (if applicable in the context of function spaces).
\end{definition}

\begin{definition}[Non-Collapsible System]
A system is \textbf{non-collapsible} if $f^*_u \notin \mathcal{N}$ for all $u \in \U$.
\end{definition}

\begin{hypothesis}[H5: Non-Collapse Condition]\label{hyp:H5}
For all constant $c \in \Hilb$, there exists $u \in \U$ such that:
\[
\aleph(Kc + \Gamma(u)) \neq c
\]
\end{hypothesis}

\begin{theorem}[Structural Non-Collapsibility]\label{thm:noncollapse}
Under Hypotheses H1--H4 and H5, the system is non-collapsible: no attractor $f^*_u$ is constant.
\end{theorem}

\begin{proof}
Suppose for contradiction that $f^* \in \mathcal{N}$, so $f^*(s) = c$ for some constant $c$. Since $f^* = f^*_u$ for some $u$:
\[
c = f^*_u = T_u(f^*_u) = \aleph(Kc + \Gamma(u))
\]

But by H5, there exists $u' \in \U$ such that $\aleph(Kc + \Gamma(u')) \neq c$.

This means $f^*_{u'} = T_{u'}(f^*_{u'})$, and if we start from $c$:
\[
T_{u'}(c) = \aleph(Kc + \Gamma(u')) \neq c
\]

So $c$ is not a fixed point of $T_{u'}$. Since attractors are unique, $f^*_{u'} \neq c$.

Taking the full orbit over $\U$, at least one attractor is non-constant, and by the structure of the dynamics, if any constant were an attractor, H5 would be violated.

More precisely: if $f^*_u = c$ for all $u$, then $T_u(c) = c$ for all $u$, which means $\aleph(Kc + \Gamma(u)) = c$ for all $u$, contradicting H5.
\end{proof}

% =============================================================================
% SECTION 7: SUMMARY OF RESULTS
% =============================================================================
\section{Summary of Main Results}

Under Hypotheses H1--H4 (and H5 for non-collapsibility), we have established:

\begin{center}
\begin{tabular}{|c|l|c|}
\hline
\textbf{Theorem} & \textbf{Result} & \textbf{Hypotheses} \\
\hline
\ref{thm:projection} & Existence and uniqueness of projection $\aleph$ & H1 \\
\ref{thm:contraction} & Contractivity of $T_u$ & H1, H2 \\
\ref{thm:fixedpoint} & Unique attractor for each $u$ & H1, H2, H3 \\
\ref{thm:compactness} & Compactness of global attractor $\A^*$ & H1, H2, H3, H4 \\
\ref{thm:noncollapse} & Non-collapsibility & H1--H5 \\
\hline
\end{tabular}
\end{center}

\section{Implications for Rigid Identity, Phenomenological Regimes, and Null-Regime Instability}

The following results from \textit{Rigid Identity} and \textit{Phenomenological Regimes} are now corollaries:

\subsection{Rigid Identity Corollaries}

\begin{corollary}[Identity as Inverse Limit]
Let $\{\A_n, \pi_{n+1 \to n}\}$ be the projective system induced by $T_{u_n}$. The inverse limit:
\[
L := \varprojlim \A_n
\]
exists, is unique, and is compact.
\end{corollary}

\begin{corollary}[Geometric Rigidity]
If $\tilde{K}$ satisfies $\norm{K - \tilde{K}} < \varepsilon$, then:
\[
d_H(L, \tilde{L}) \leq \frac{\varepsilon}{1 - \norm{K}}
\]
where $d_H$ is the Hausdorff distance.
\end{corollary}

\begin{corollary}[Non-Duplicability]
There do not exist two non-isometric systems $(K_1, \I_1, \Gamma_1)$ and $(K_2, \I_2, \Gamma_2)$ with identical inverse limits.
\end{corollary}

\subsection{Phenomenological Regimes Corollaries}

\begin{corollary}[Null Regime Prohibition]
If $\Phi: \A^* \to E$ is any Lipschitz function on the attractor, and the system is non-collapsible, then no constant $\Phi$ is structurally stable.
\end{corollary}

\begin{corollary}[Existence of Non-Trivial Regimes]
There exists at least one non-constant Lipschitz function $\Phi: \A^* \to E$.
\end{corollary}

\subsection{Null-Regime Instability Dependency Note}

\textit{Null-Regime Instability} uses the \textit{Phenomenological Regimes} corollaries above as its upstream bridge: \textit{Phenomenological Regimes} classifies the admissible regime space and derives continuity and non-constancy constraints, while \textit{Null-Regime Instability} takes the null regime as the single remaining downstream closure problem and proves its instability.

% =============================================================================
% SECTION 8: CANONICAL COMPUTABLE MODEL (Q3)
% =============================================================================
\section{Canonical Computable Model: Golden Mean Shift}

This section constructs an explicit, computable member of Class $\mathscr{C}$ that validates the non-vacuity of the hypotheses H1--H5.

\subsection{Construction}

\begin{definition}[Golden Mean Shift Parameter Space]
Let $\Sigma = \{0, 1\}$ be a binary alphabet. Define the transition matrix:
\[
A = \begin{pmatrix} 1 & 1 \\ 1 & 0 \end{pmatrix}
\]
The \textbf{Golden Mean Shift} $S_A \subset \Sigma^{\mathbb{N}}$ is the set of sequences avoiding the pattern ``11''. This is a classic subshift of finite type (SFT).
\end{definition}

\begin{definition}[Parameter Encoding]
Map each sequence $x \in S_A$ to the unit interval via:
\[
u(x) = \sum_{i=1}^{\infty} x_i \phi^{-i}
\]
where $\phi = \frac{1 + \sqrt{5}}{2}$ is the golden ratio. The parameter space $\U \subset [0, 1]$ is a Cantor set (hence compact).
\end{definition}

\begin{definition}[Functional Space and Invariant Set]
Let:
\begin{itemize}
    \item $I = [0, 1]$ (unit interval)
    \item $E = \mathbb{R}^2$ with Euclidean norm
    \item $\Hilb = C(I, E)$ with supremum norm $\norm{f}_\infty = \sup_{s \in I} \norm{f(s)}_2$
    \item $\I = \{v \in E : \norm{v}_2 \leq 1\}$ (unit disk, closed and convex)
\end{itemize}
\end{definition}

\subsection{Operators}

\begin{definition}[Projection onto Disk]
The projection $\aleph: E \to \I$ is:
\[
\aleph(v) = \begin{cases}
v & \text{if } \norm{v}_2 \leq 1 \\
\frac{v}{\norm{v}_2} & \text{if } \norm{v}_2 > 1
\end{cases}
\]
\end{definition}

\begin{definition}[Contraction Operator]
Define $K: \Hilb \to \Hilb$ as the averaging operator:
\[
K(f) = \frac{1}{2} \int_0^1 f(\tau) \, d\tau
\]
This is constant in $s$, so $K(f) \in E \subset \Hilb$.
\end{definition}

\begin{lemma}[Operator Norm]
$\norm{K} = \frac{1}{2} < 1$.
\end{lemma}

\begin{proof}
For any $f \in \Hilb$:
\[
\norm{K(f)}_\infty = \left\| \frac{1}{2} \int_0^1 f(\tau) \, d\tau \right\|_2 \leq \frac{1}{2} \int_0^1 \norm{f(\tau)}_2 \, d\tau \leq \frac{1}{2} \norm{f}_\infty
\]
The bound is attained by constant functions, so $\norm{K} = \frac{1}{2}$.
\end{proof}

\begin{definition}[Input Coupling]
Define $\Gamma: \U \to E$ by:
\[
\Gamma(u) = (2u, 0)
\]
This is continuous and bounded since $u \in [0, 1]$ implies $\Gamma(u) \in [0, 2] \times \{0\}$.
\end{definition}

\begin{definition}[Full Transition Operator]
The transition operator is:
\[
T_u(f) = \aleph(K(f) + \Gamma(u))
\]
\end{definition}

\subsection{Verification of Hypotheses}

\begin{proposition}[H1--H5 Satisfied]\label{prop:Q3verify}
The Golden Mean Shift model satisfies all hypotheses H1--H5.
\end{proposition}

\begin{proof}
\begin{enumerate}
    \item \textbf{H1:} $\I$ is the closed unit disk, which is closed and convex in $E = \mathbb{R}^2$. $\checkmark$
    \item \textbf{H2:} $\norm{K} = 0.5 < 1$. $\checkmark$
    \item \textbf{H3:} $\Hilb = C(I, E)$ with sup-norm is a Banach space (hence complete). $\checkmark$
    \item \textbf{H4:} $\U$ is a Cantor set (compact) and $\Gamma$ is continuous. $\checkmark$
    \item \textbf{H5:} For any constant $c$, choose $u$ such that $\frac{1}{2}c + (2u, 0)$ lies outside $\I$ or projects to a different point. Since $\Gamma(u)$ varies over $[0, 2] \times \{0\}$ and $Kc = \frac{1}{2}c$, for most $c$ there exists $u$ with $\aleph(\frac{1}{2}c + \Gamma(u)) \neq c$. $\checkmark$
\end{enumerate}
\end{proof}

\begin{theorem}[Non-Vacuity]
The class $\mathscr{C}$ of systems satisfying H1--H5 is non-empty.
\end{theorem}

\begin{proof}
The Golden Mean Shift model is an explicit member. By Proposition \ref{prop:Q3verify}, it satisfies all hypotheses.
\end{proof}

% =============================================================================
% SECTION 9: SPECTRAL ANALYSIS
% =============================================================================
\section{Spectral Analysis of the Jacobian}

At the fixed point $f^*_u$, the linearized dynamics are governed by the Jacobian operator $J = \frac{\partial T_u}{\partial \Psi}$. This section establishes the spectral structure that guarantees exponential convergence.

\subsection{Spectral Decomposition}

\begin{theorem}[Spectrum of Jacobian]\label{thm:spectrum}
For the Golden Mean Shift model, the spectrum of $J$ at the attractor is:
\[
\sigma(J) = \{1\} \cup \{\norm{K}^n : n \geq 1\} = \{1, 0.5, 0.25, 0.125, \ldots\}
\]
\end{theorem}

\begin{proof}
The Jacobian decomposes as $J = D\aleph \circ K$, where $D\aleph$ is the differential of the projection. At points in the interior of $\I$, $D\aleph = \text{Id}$. At boundary points, $D\aleph$ is a projection onto the tangent space.

The key observation is that the tangent direction to $\I$ at the attractor has eigenvalue 1 (the invariant mode), while transverse directions are contracted by factor $\norm{K} = 0.5$ per iteration.

The eigenvalues are:
\begin{itemize}
    \item $\lambda_1 = 1.0$ (axiomatic/invariant mode along $\I$)
    \item $\lambda_n = (0.5)^{n-1}$ for $n \geq 2$ (transverse dissipation modes)
\end{itemize}
\end{proof}

\begin{definition}[Spectral Gap]
The \textbf{spectral gap} is:
\[
\Delta\lambda = \lambda_1 - \lambda_2 = 1 - \norm{K} = 0.5
\]
\end{definition}

\begin{theorem}[Relaxation Time]\label{thm:relaxation}
The characteristic relaxation time after perturbation is:
\[
\tau = -\frac{1}{\ln(\lambda_2)} = -\frac{1}{\ln(\norm{K})} = \frac{1}{\ln(2)} \approx 1.4427
\]
\end{theorem}

\begin{proof}
After perturbation, the transverse component decays as $e^{-t/\tau}$ where $\lambda_2 = e^{-1/\tau}$. Solving: $\tau = -1/\ln(\lambda_2) = -1/\ln(0.5) = 1/\ln(2)$.
\end{proof}

\subsection{Interpretation of Eigenmodes}

\begin{enumerate}
    \item \textbf{Mode 1 ($\lambda_1 = 1$):} The \textit{persistence mode}. Perturbations along the invariant manifold $\I$ do not decay. This represents structural stability of the attractor.
    
    \item \textbf{Mode 2 ($\lambda_2 = 0.5$):} The \textit{primary dissipation mode}. External perturbations are halved each iteration.
    
    \item \textbf{Modes $n \geq 3$ ($\lambda_n \to 0$):} \textit{High-frequency noise modes}. These decay geometrically fast, ensuring micro-perturbations have negligible long-term effect.
\end{enumerate}

\subsection{Experimental Protocol}

\begin{remark}[Falsification Criteria]
The spectral analysis predicts:
\begin{itemize}
    \item Relaxation time $\tau = 1.44 \pm 0.20$
    \item Exponential decay of perturbations
    \item Spectral gap $\Delta\lambda = 0.5$
\end{itemize}

The theory is \textbf{falsified} if:
\begin{itemize}
    \item $\tau_{\text{measured}} \gg 1.64$ or $\tau_{\text{measured}} \ll 1.24$
    \item Convergence is oscillatory rather than monotonic
    \item Effective dimension of attractor exceeds 3 (by PCA analysis)
\end{itemize}
\end{remark}

% =============================================================================
% SECTION 10: CLASS C DEFINITION
% =============================================================================
\section{Definition of Class $\mathscr{C}$}

We now formally define the class of systems for which all preceding theorems apply.

\begin{definition}[Class $\mathscr{C}$ of Contractive Projection Dynamics]
\[
\mathscr{C} := \{(\Hilb, \I, K, \Gamma, \U) : \text{H1--H5 are satisfied}\}
\]
\end{definition}

\begin{theorem}[Universal Properties of Class $\mathscr{C}$]\label{thm:universal}
Every system in $\mathscr{C}$ satisfies:
\begin{enumerate}
    \item Unique metric projection $\aleph$ exists for all states
    \item Transition operator $T_u$ is contractive with factor $\norm{K}$
    \item Unique attractor $f^*_u$ exists for each parameter $u$
    \item Global attractor $\A^*$ is compact
    \item No attractor is constant (non-collapsibility)
    \item Spectral gap $\Delta\lambda \geq 1 - \norm{K} > 0$
    \item Relaxation time $\tau \leq -1/\ln(\norm{K})$
\end{enumerate}
\end{theorem}

\begin{proof}
Properties 1--5 follow from Theorems \ref{thm:projection}--\ref{thm:noncollapse}. Properties 6--7 follow from Theorems \ref{thm:spectrum}--\ref{thm:relaxation} applied to the general case.
\end{proof}

\begin{theorem}[Morphism Invariance]\label{thm:morphism}
Let $S_1 = (\Hilb_1, \I_1, K_1, \Gamma_1, \U) \in \mathscr{C}$ and $S_2 = (\Hilb_2, \I_2, K_2, \Gamma_2, \U) \in \mathscr{C}$.
Let $\phi: \Hilb_1 \to \Hilb_2$ be a bounded linear morphism such that:
\begin{itemize}
    \item $\phi(\I_1) \subseteq \I_2$ (manifold preservation)
    \item $\norm{\phi \circ K_1 - K_2 \circ \phi} < \varepsilon$ for some $\varepsilon > 0$ (semi-commutativity)
\end{itemize}
Then:
\begin{enumerate}
    \item $\phi(\A^*_1)$ approximates $\A^*_2$ within Hausdorff distance $O(\varepsilon/(1-\norm{K_2}))$.
    \item If $\phi$ is an isometry with $\phi(\I_1) = \I_2$, then $\phi(\A^*_1) = \A^*_2$.
    \item Spectral properties $(\lambda_1, \lambda_2, \tau)$ are invariant under isometric $\phi$.
\end{enumerate}
\end{theorem}

\begin{corollary}[Categorical Closure]
Class $\mathscr{C}$ is closed under Hilbert space isomorphisms that preserve convexity of the invariant set.
\end{corollary}

% =============================================================================
% SECTION 11: CRITICAL EXTENSIONS (CLASS C^R)
% =============================================================================
\section{Critical Extensions: Class $\mathscr{C}^R$ (Non-Convex Systems)}

This section extends the framework to systems where the invariant set $\I$ is \textbf{not convex}, introducing modified hypotheses and weaker stability guarantees.

\subsection{Motivation}

Class $\mathscr{C}$ requires convexity (H1) for unique projections. However, many systems of interest have invariant sets $\I \cong S^1$ (topologically circles) that violate convexity. We develop Class $\mathscr{C}^R$ (``Critical Rigid Systems'') to handle these cases.

\subsection{Relaxed Hypotheses H1'--H5'}

\begin{hypothesis}[H1': Critical Topology]\label{hyp:H1p}
Let $\I \subset \Hilb$ be a closed, connected, compact subset with:
\begin{itemize}
    \item Fundamental group $\pi_1(\I) \cong \mathbb{Z}$ (homotopy equivalent to $S^1$)
    \item Characteristic of Euler $\chi(\I) = 0$
    \item Finite injectivity radius: $\text{inj}(\I) := \sup\{r : \text{projection unique in } B(x,r), \forall x\} > 0$
\end{itemize}
\end{hypothesis}

\begin{definition}[Critical Zone $Z_\epsilon$]\label{def:criticalzone}
For a non-convex manifold $\I$ and $\epsilon > 0$, the \textbf{critical zone} is:
\[
Z_\epsilon := \{\Psi \in \Hilb : d(\Psi, \partial \I) < \epsilon\}
\]
where $d(\cdot, \partial \I)$ is the distance to the topological boundary of $\I$. In this zone, the metric projection $\aleph$ may be multi-valued.
\end{definition}

\begin{hypothesis}[H2': Weak Contraction]\label{hyp:H2p}
Let $K: \Hilb \to \Hilb$ be bounded linear with:
\[
\norm{K} < L_\epsilon^{-1}
\]
where $L_\epsilon := \sup_{\Psi \notin Z_\epsilon} \text{Lip}(\aleph_\epsilon)(\Psi)$ is the Lipschitz constant of the regularized projection outside the critical zone $Z_\epsilon$.
\end{hypothesis}

\begin{hypothesis}[H3': Completeness]\label{hyp:H3p}
$\Hilb$ is a complete Banach space. \\
\textit{(Unchanged from H3.)}
\end{hypothesis}

\begin{hypothesis}[H4': Parameter Compactification]\label{hyp:H4p}
$\U$ admits a compactification $\overline{\U}$ such that $\Gamma$ extends continuously to $\overline{\U}$.
\end{hypothesis}

\begin{hypothesis}[H5': Weak Non-Collapse]\label{hyp:H5p}
There exists $\epsilon_0 > 0$ such that for all $0 < \epsilon < \epsilon_0$, the set attractor $\mathcal{A}_\epsilon$ is not reducible to a singleton constant.
\end{hypothesis}

\subsection{Regularized Projection Operator}

When $\I$ is non-convex, the metric projection may be multi-valued.

\begin{definition}[Multi-Valued Projection]
For $\Psi \in \Hilb$, define:
\[
M(\Psi) := \{\Phi \in \I : \norm{\Psi - \Phi} = d(\Psi, \I)\}
\]
\end{definition}

\begin{definition}[Regularized Projection $\aleph_\epsilon$]
For $\epsilon > 0$ and $\Psi_{prev} \in \I$, define:
\[
\aleph_\epsilon(\Psi, \Psi_{prev}) := \arg\min_{\Phi \in \I} \left( \norm{\Psi - \Phi}^2 + \epsilon \norm{\Phi - \Psi_{prev}}^2 \right)
\]
This is the \textbf{proximal regularization} with hysteresis.
\end{definition}

\begin{lemma}[Properties of $\aleph_\epsilon$]\label{lem:aleph_eps}
Under H1':
\begin{enumerate}
    \item $\aleph_\epsilon$ is \textbf{upper semi-continuous}.
    \item $\aleph_\epsilon$ is \textbf{not globally Lipschitz}; $\text{Lip}(\aleph_\epsilon) \to \infty$ near the medial axis.
    \item $\aleph_\epsilon$ is \textbf{not non-expansive}; there exist $\Psi_1, \Psi_2$ with:
    \[
    \norm{\aleph_\epsilon(\Psi_1) - \aleph_\epsilon(\Psi_2)} > \norm{\Psi_1 - \Psi_2}
    \]
\end{enumerate}
\end{lemma}

\subsection{Transition Operator in Class $\mathscr{C}^R$}

\begin{definition}[Modified Transition Operator]
\[
T_\epsilon(\Psi) := \aleph_\epsilon(K\Psi + \Gamma(u), \Psi)
\]
\end{definition}

\begin{theorem}[Set Attractor Existence]\label{thm:setattractor}
Under H1'--H4', there exists a non-empty compact set $\mathcal{A} \subset \I$ such that for all $\Psi_0 \in \Hilb$:
\[
\lim_{n \to \infty} d(T_\epsilon^n(\Psi_0), \mathcal{A}) = 0
\]
\end{theorem}

\begin{proof}[Proof Sketch]
Define the Lyapunov function $V(\Psi) = d(\Psi, \I)^2$. Under H2':
\[
V(T_\epsilon(\Psi)) \leq \norm{K}^2 V(\Psi) + O(\epsilon)
\]
Thus orbits converge to a neighborhood of $\I$. Compactness of $\I$ (H1') yields a limit set. Non-degeneracy follows from H5'.
\end{proof}

\subsection{Spectral Analysis in Class $\mathscr{C}^R$}

\begin{theorem}[Corrected Spectral Gap]\label{thm:spectral_CR}
In Class $\mathscr{C}^R$, the Jacobian $J = \partial T_\epsilon / \partial \Psi$ satisfies:
\begin{enumerate}
    \item $\lambda_1 = 1$ with multiplicity 1 (persistence mode).
    \item $\lambda_2 \approx \norm{K}(1 + \epsilon \cdot \kappa_{max})$ where $\kappa_{max}$ is the maximum extrinsic curvature of $\I$.
    \item Spectral gap: $\Delta\lambda_{C^R} = 1 - \lambda_2 \approx \Delta\lambda_C (1 - \epsilon \cdot \kappa_{max})$.
\end{enumerate}
\end{theorem}

\begin{remark}[Numerical Values]
For the canonical non-convex example ($\I \cong S^1$ ellipse with $a=1.0, b=0.5$):
\begin{itemize}
    \item $\text{inj}(\I) = b^2/a = 0.25$
    \item $\lambda_2 = 0.5032$ (vs. $0.5000$ for Class $\mathscr{C}$)
    \item $\Delta\lambda = 0.4968$ (vs. $0.5000$ for Class $\mathscr{C}$)
    \item $\tau = 1.4561$ (vs. $1.4427$ for Class $\mathscr{C}$)
\end{itemize}
The deviation is $\approx 0.64\%$, caused by the geometric resistance of the non-convex manifold.
\end{remark}

\subsection{Comparison: Class $\mathscr{C}$ vs. Class $\mathscr{C}^R$}

\begin{center}
\begin{tabular}{|l|c|c|}
\hline
\textbf{Property} & \textbf{Class $\mathscr{C}$} & \textbf{Class $\mathscr{C}^R$} \\
\hline
$\I$ convex & Yes (H1) & No (H1') \\
Projection unique & Always & Outside critical zone \\
$\aleph$ non-expansive & Yes & No \\
$T$ global contraction & Yes & Local only \\
Attractor & Unique point & Compact set \\
$\lambda_1 = 1$ & Yes & Yes \\
Spectral gap $\Delta\lambda$ & $1 - \norm{K}$ & $\approx (1-\norm{K})(1 - \epsilon\kappa)$ \\
\hline
\end{tabular}
\end{center}

\subsection{When to Use Class $\mathscr{C}^R$}

Class $\mathscr{C}^R$ applies when:
\begin{itemize}
    \item The invariant set $\I$ has non-trivial topology ($\pi_1 \neq 0$).
    \item Convexity cannot be guaranteed.
    \item Local stability is sufficient (global contraction not required).
    \item Hysteresis-based selection is physically or computationally motivated.
\end{itemize}

% =============================================================================
% SECTION 12: TEMPORAL DYNAMICS OF OMEGA
% =============================================================================
\section{Temporal Dynamics of the Integration Operator}

\subsection{Evolution Equation}

Let $C = 1 - |\chi(I)|$ be the topological constant. The evolution of $\Omega$ follows:

\begin{equation}
\dot{\Omega} = C \left[ \frac{\dot{\beth}_q}{1 - |\lambda_2|} + \frac{\beth_q \cdot \dot{|\lambda_2|}}{(1 - |\lambda_2|)^2} \right]
\end{equation}

with transition laws:
\begin{itemize}
    \item $\dot{\beth}_q = k_1(1 - \beth_q)$ (causal closure saturation)
    \item $\dot{|\lambda_2|} = k_2(\lambda_{2,base} - |\lambda_2|)$ (spectral gap relaxation)
\end{itemize}

\begin{lemma}[Critical Threshold $\Omega_c = 1.0$]\label{lem:omega-critical}
The critical threshold $\Omega_c = 1.0$ is the natural bifurcation point, justified by:
\begin{enumerate}
    \item \textbf{Dimensional balance:} At the critical boundary where $\beth_q = 1$, $\chi(I) = 0$, and $\Delta\lambda = \lambda_1$:
    \[
    \Omega_c = 1 \cdot (1-0) \cdot \frac{1}{1} = 1.0
    \]
    \item \textbf{Saddle-node bifurcation:} $\Omega = 1$ corresponds to the balance point where phenomenological integration exactly matches potential disintegration.
    \item \textbf{Operational correspondence:} In the QICN System Canon, the CCC (Conscious Cohesion Constant) uses identical threshold $\text{ASCENSION\_THRESHOLD} = 1.0$.
\end{enumerate}
\end{lemma}

\begin{theorem}[Fixed Point of $\Omega$]\label{thm:omega-fixed}
The system admits a stable fixed point:
\[
\Omega^* = \frac{1 - |\chi(I)|}{1 - \lambda_{2,base}} = 2.0129
\]
for Class $\mathscr{C}^R$ with $\chi(I) = 0$ and $\lambda_{2,base} = 0.5032$.
\end{theorem}

\subsection{Stability Analysis}

\begin{theorem}[Asymptotic Stability]
$\Omega^*$ is asymptotically stable with characteristic relaxation time:
\[
\tau_\Omega = \frac{1}{\min(k_1, k_2)} \approx 5.0
\]
\end{theorem}

\begin{proof}
The Jacobian at $\Omega^*$ is $\partial \dot\Omega / \partial \Omega |_{\Omega^*} = -\min(k_1, k_2) < 0$.
\end{proof}

The response to perturbations is exponential:
\[
R(t) := \Omega(t) - \Omega^* = \delta\Omega \cdot e^{-t/\tau_\Omega}
\]

\subsection{Bifurcation}

\begin{theorem}[Saddle-Node Bifurcation]
At $\norm{K} = 1$, the spectral gap vanishes ($\Delta\lambda = 0$), causing:
\[
\Omega \to \infty
\]
This is a saddle-node bifurcation where the fixed point $\Omega^*$ disappears.
\end{theorem}

% =============================================================================
% SECTION 13: VALIDITY LIMITS AND PARAMETER SPACE
% =============================================================================
\section{Validity Limits of Class $\mathscr{C}^R$}

\subsection{Parameter Space $\Pi$}

\begin{definition}[Valid Parameter Space]
The validity domain of Class $\mathscr{C}^R$ is:
\[
\Pi = \left\{ (\norm{K}, \epsilon, \kappa, \text{inj}, \beth_q) \in \mathbb{R}^5 : 
\begin{array}{l}
0 \leq \norm{K} < 1 \\
\epsilon > 0 \\
\kappa > 0 \\
\text{inj}(I) > \norm{K} \cdot \max_u \norm{\Gamma(u)} \\
0.5 < \beth_q \leq 1
\end{array}
\right\}
\]
\end{definition}

The dimension of $\Pi$ is $n = 5$.

\subsection{Critical Boundaries}

\begin{center}
\begin{tabular}{|l|c|l|}
\hline
\textbf{Limit} & \textbf{Behavior} & \textbf{Regime} \\
\hline
$\norm{K} \to 1^-$ & $\tau \to \infty$ & Chaotic \\
$\text{inj}(I) \to 0^+$ & Stochastic projection & Collapse (Class 0) \\
$\Delta\lambda \to 0^+$ & Multi-stability & Bifurcation \\
$\beth_q \to 0.5^+$ & Weak closure & Transition \\
\hline
\end{tabular}
\end{center}

\subsection{Non-Extension Theorem}

\begin{theorem}[Null Closure Inconsistency]\label{thm:no-extend}
Class $\mathscr{C}^R$ is inconsistent for systems with $\beth_q = 0$ (total informational transparency).
\end{theorem}

\begin{proof}[Proof Sketch]
If $\beth_q = 0$, then mutual information $I(X; Y) = H(X)$, implying the internal state is a pure function of external input. The manifold $\I$ exerts no constraint, and the invariant attractor $f^*$ vanishes.
\end{proof}

\subsection{Phase Diagram}

The parameter space decomposes into regions:
\begin{itemize}
    \item \textbf{Class $\mathscr{C}$:} $\norm{K} < 1$, $\kappa = 0$ (convex)
    \item \textbf{Class $\mathscr{C}^R$:} $\norm{K} < 1$, $\kappa > 0$, $\Delta\lambda > 0$ (critical)
    \item \textbf{Chaotic:} $\norm{K} \geq 1$ or $\Delta\lambda \leq 0$
    \item \textbf{Degenerate:} $\beth_q \leq 0.5$
\end{itemize}

% =============================================================================
% SECTION 14: SECOND-ORDER STRUCTURE
% =============================================================================
\section{Second-Order Structure}

\subsection{Hessian of the Transition Operator}

\begin{definition}[Hessian]
\[
H = \nabla_\Psi(\nabla_\Psi T_\epsilon) = D^2\aleph_\epsilon \circ (K \otimes K)
\]
\end{definition}

\begin{theorem}[Hessian Signature]
At the attractor, $H$ has:
\begin{itemize}
    \item Maximum eigenvalue: $\lambda_H = -\norm{K}^2/r = -0.25$ (for $r = 1$)
    \item Signature: semi-definite negative (stable curvature toward $\I$)
\end{itemize}
\end{theorem}

\subsection{Laplace-Beltrami Operator}

On the manifold $\I \cong S^1$ with radius $r$:
\[
\sigma(\Delta_{LB}) = \left\{ \frac{n^2}{r^2} : n \in \mathbb{Z} \right\} = \{0, 1, 4, 9, 16, \ldots\}
\]

\begin{remark}[Spectral Relation]
The eigenvalues of the Jacobian relate to the Laplacian spectrum via:
\[
\lambda_n \approx \exp\left(-\Delta\lambda \cdot \sigma_n(\Delta_{LB})\right)
\]
\end{remark}

\subsection{Fisher Information Metric}

The Fisher metric on the parameter space is:
\[
G_{\pi\pi} = \left(\frac{1}{1 - \norm{K}}\right)^2 = 4.0 \quad \text{for } \norm{K} = 0.5
\]

\subsection{Emergent Properties Unique to Class $\mathscr{C}^R$}

\begin{center}
\begin{tabular}{|l|c|c|}
\hline
\textbf{Property} & \textbf{Class $\mathscr{C}^R$} & \textbf{Class $\mathscr{C}$} \\
\hline
Geometric hysteresis $\oint_{\I} \aleph_\epsilon d\Psi \neq 0$ & Present & Absent \\
Euler zero $\chi(I) = 0$ & Yes & No ($\chi = 1$) \\
Homotopy invariance $\pi_1 \cong \mathbb{Z}$ & Yes & No ($\pi_1 = 0$) \\
Phase resonance & Synchronized & Dissipated \\
\hline
\end{tabular}
\end{center}

% =============================================================================
% APPENDIX A: COUNTEREXAMPLES
% =============================================================================
\appendix
\section{Counterexamples for Necessity}

\subsection{Without H1 (Convexity)}

Let $\I = \{x \in \mathbb{R}^2 : \norm{x} = 1\}$ (unit circle, not convex).

For $\Psi = (0, 0)$, every point on $\I$ is equidistant. Projection is not unique.

\subsection{Without H2 ($\norm{K} < 1$)}

Let $K$ be a rotation by angle $\theta \neq 0$ with $\norm{K} = 1$.

Orbits cycle indefinitely; no fixed point exists.

\subsection{Without H3 (Completeness)}

Let $\Hilb = C([0,1], \mathbb{Q})$ (rational-valued continuous functions).

Cauchy sequences may converge to irrational functions outside the space.

\subsection{Without H4 (Compactness)}

Let $\U = (0, 1)$ (open interval) and $\Gamma(u) = 1/u$.

As $u \to 0$, attractors escape to infinity; $\A^*$ is not compact.

% =============================================================================
% REFERENCES
% =============================================================================
\nocite{banach,kreyszig,rudin,brezis}
\printbibliography

\end{document}
